\section{Related Work}
\label{sec:related}

Prior work runs along four threads: the measurement of geographic factual bias,
multilingual and cultural extensions, mitigation attempts, and LLM use in policy
and health decision contexts.

\subsection{Geographic factual bias and multilingual extensions}

\citet{manvi2024llm} established that frontier LLMs carry systematic biases
against regions of lower socioeconomic conditions, with correlations up to 0.70
between model-generated ratings and a proxy for those conditions; the outcome
space is subjective ratings, which leaves open whether the same bias appears in
objective recall of regulatory facts. \citet{moayeri2024worldbench} quantified
factual recall against World Bank indicators across 20 LLMs, documenting
1.5$\times$ higher absolute relative error for Sub-Saharan Africa than for North
America; the design is bounded by the macro-economic indicator space, while an
environmental manager asks narrower, source-specific questions whose ground
truth lives in environmental councils and monitoring agencies. \citet{mirza2024global}
showed that newer GPT-4 versions do not monotonically improve and that the
privileging of Global North statements persists.

The multilingual probing literature established early that what a model
recalls depends on the language of the query. \citet{jiang2020xfactr} probed
cloze-style facts in 23 languages and found accuracy below 15\% even in English
and Spanish and below 5\% in Marathi and Yoruba, with half of the correctly
recalled facts recalled in one language only; \citet{kassner2021mlama}
translated the same benchmarks into 53 languages and found mBERT below 60\% of
its English performance in 32 of them, a bias toward entities of the query
language, and no clear relation between the size of a language's Wikipedia and
its recall. Those results concern masked encoders of 2020; the question they
raise, whether recall in a language tracks the corpus of that language, is the
one H2 and H4 put to generative models. \citet{myung2024blend} covered 16
countries and 13 languages in everyday cultural knowledge, finding that LLMs perform better in the
native languages of high-resource cultures and worse in those of low-resource
ones, the finding that motivates H2 and connects to the resource hierarchy of
\citet{joshi2020state}. Regional benchmarks have proliferated:
BRoverbs~\citep{broverbs2025}, TiEBe~\citep{almeida2025tiebe},
ALBA~\citep{alba2026} and AMALIA~\citep{amalia2026} document concrete failures in
culturally grounded tasks but do not address applied policy use. Work on
Portuguese corpus construction is explicit that the effort has centred on
English: \citet{almeida2025portuguese} assemble a 120B-token Portuguese corpus
because language-specific pipelines are needed to reach industrial quality. That
asymmetry in how much text exists per language is what H4 tests.

These studies establish that the phenomenon exists and is measurable. None
isolates a single applied domain with verifiable official ground truth, and none
treats prompt framing as a manipulable factor.

\subsection{Mitigation, persona, and role conditioning}

\citet{opuszko2026unraveling} found that reasoning helps in aggregate but does
not eliminate regional disparities, and \citet{kerche2026silicon} proposed a
place-based typology situating geographic bias in a broader sociotechnical
landscape. Both are conceptual and aggregate.

Role conditioning is what the vendors instruct their own users to do. Anthropic
documents role setting in the system prompt~\citep{anthropic2026systemprompts},
OpenAI's guide opens the recommended message structure with an identity
block~\citep{openai2026promptengineering}, Google's Gemini documentation asks
that role definitions, which it calls persona, be placed in the system
instruction~\citep{google2026promptingstrategies}, xAI's standard example
initialises every conversation with a role-assigning system
prompt~\citep{xai2026chatguide}, and the pattern is catalogued in the prompting
literature~\citep{white2023prompt}. The systematic evidence on whether it helps
is negative: across four LLM families and 2{,}410 factual questions, expert
personas did not improve accuracy over a no-persona control and slightly reduced
it across 162 roles~\citep{zheng2024helpful}. An independent study with six
models on graduate-level multiple-choice questions reaches the same conclusion:
matching an expert persona to the problem type had no significant effect on
performance~\citep{mollick2025personas}. Whether persona conditioning helps or
harms factual reliability in high-stakes domains is therefore unsettled and, to
our knowledge, untested for geographic equity. We treat persona as a
pre-specified experimental factor and test both the Norm-Compliance hypothesis
(persona reduces the gap) and the Persona-Vacuum hypothesis (persona amplifies
it).

\subsection{LLMs in policy, health, and the administration}

\citet{lecoz2025policymaking} evaluated LLM policy recommendations for
homelessness across four cities (Barcelona, Johannesburg, South Bend and Macau)
and read the pattern as ``a context-blind rigidity and a potential bias based on
the published online discourse on homelessness in the regions of the world with
overrepresented data in LLMs'', against experts who tailored decisions to
locally encoded realities. That is one domain and one task type, which we
generalise to five task types spanning recall, synthesis and recommendation. The closest precedent in a
regulated domain is \citet{dahl2024legal}, who asked four public models
verifiable questions about United States federal cases and found hallucination
rates between 58\% and 88\%, highest for the lowest courts and least prominent
jurisdictions, and highest again for the oldest and the most recent cases, so
that the models' knowledge lags the doctrine in force; our recall floor and the
ladder analysis of Section~\ref{sec:results:t1} are the environmental
counterpart of those two gradients. In
health, \citet{kozlakidis2026medical} argue that hallucinations carry acute risk
in regulatory environments with limited technical capacity, where language and
cultural mismatch increases error, a risk that meets the mortality burden
documented by the Global Burden of Disease estimates~\citep{gbd2019risk}.

The deployment question is no longer hypothetical. \citet{kim2026genai}
estimates bureaucratic productivity effects of generative AI
quasi-experimentally, and \citet{robinson2026opensource} documents that choosing
which model an agency runs is now an explicit procurement decision between
open-weight and vendor-hosted options. Our model-level findings speak to that
decision: open-weight models as a group trail the closed ones by
$\nhfive$~points on the composite even though one open-weight model
(DeepSeek-V3) leads the roster, and the regionally tuned Portuguese model is the
weakest of the fourteen. Two further strands frame why a wrong answer matters
administratively. \citet{choroszewicz2026crumple} shows that generative-AI
support tools position frontline practitioners as ``moral crumple zones'',
absorbing responsibility for failures originating upstream, and
\citet{cordella2024regulating}, analysing the first binding measure a regulator
imposed on a generative-AI service (Italy's intervention on ChatGPT), argue that
technology-neutral frameworks overlook what is specific to these systems and
that regulators must engage their technical workings directly. On reception,
\citet{aoki2024explainable} test whether the type of explanation changes
perceived accuracy, fairness and trustworthiness among those affected.

None of this establishes whether the underlying answer is correct, or whether
correctness is distributed evenly across the states that rely on it. Perceived
trustworthiness and measured accuracy are different quantities, and a system can
be well explained, well governed, well received, and wrong. That is the gap this
study addresses, in a domain where the correct answer is published in an
official register and can be checked.

\subsection{Positioning of the present work}

The design differs from prior work in four ways: a single high-stakes applied
domain with verifiable official ground truth rather than numeric trivia or
cultural knowledge; persona as a pre-specified within-subject factor rather than
an uncontrolled framing choice; theory-driven country sampling on independent
axes (the UNCTAD developed/developing split, Joshi linguistic resource,
development) rather than regional convenience; and an explicit
corpus-representation mechanism test with sensitivity analysis. No prior study
combines a regulated public-health domain, a persona manipulation, and a
pre-specified mechanism test.
